\chapter[Geração dos relatórios]{Geração dos relatórios} \label{capitulo12}

A geração dos relatórios foi implementada como uma etapa posterior ao registro das aplicações das avaliações, permitindo que o profissional organize e documente suas considerações a partir das informações obtidas durante o processo avaliativo. Essa funcionalidade está relacionada ao objetivo do sistema de apoiar a sistematização das informações e a elaboração de devolutivas individualizadas, mantendo a interpretação dos dados sob responsabilidade do profissional.

No modelo de dados, os relatórios são representados pela entidade \textit{relatorios}, que possui uma relação de um para um com \textit{aplicacoes}. Cada relatório está, portanto, associado a uma aplicação específica, permitindo identificar o instrumento utilizado, o aprendente relacionado e o contexto em que os dados foram registrados. A entidade também mantém a referência ao profissional responsável pela elaboração do relatório.

Os principais campos destinados ao conteúdo do relatório são \textit{conclusao} e \textit{observacoesProfissional}. O campo \textit{conclusao} permite registrar a síntese elaborada pelo profissional a partir do processo avaliativo, enquanto \textit{observacoesProfissional} possibilita acrescentar informações complementares consideradas relevantes. A estrutura também possui campos de controle de criação e atualização, permitindo acompanhar o registro e suas alterações.

A associação direta entre relatório e aplicação contribui para a rastreabilidade das informações. Como cada aplicação identifica o aprendente, o instrumento utilizado, o profissional responsável, a data e as respostas registradas, o relatório pode ser contextualizado dentro do atendimento correspondente. Essa organização evita que a devolutiva seja tratada como um registro isolado e permite relacioná-la às evidências coletadas durante a aplicação da avaliação.

O sistema também mantém a separação entre os dados obtidos por meio dos instrumentos e a interpretação profissional desses dados. As respostas das questões são armazenadas na entidade \textit{respostas}, enquanto as conclusões e observações são registradas posteriormente em \textit{relatorios}. Essa separação é importante porque os resultados de uma avaliação não são convertidos automaticamente pelo sistema em diagnósticos ou conclusões psicopedagógicas. A análise e a interpretação permanecem como atribuições do profissional responsável pelo atendimento.

A funcionalidade de relatórios também está relacionada à necessidade de produzir registros individualizados, identificada durante a definição dos requisitos do sistema. A proposta apresentada no TCC considera que as conclusões psicopedagógicas são contextualizadas e envolvem interpretação profissional, tornando necessário um mecanismo que permita registrar informações qualitativas juntamente aos dados estruturados obtidos durante as avaliações.

Dessa forma, o módulo de relatórios complementa o fluxo de avaliação implementado no sistema, estabelecendo uma sequência entre a aplicação do instrumento, o registro das respostas e a elaboração da devolutiva profissional. A estrutura desenvolvida permite centralizar essas informações e manter sua associação ao histórico do aprendente, contribuindo para a organização e a continuidade dos registros avaliativos.